\begin{statement}{C.2.MH.2}{Proved}{Euler number from a ruled normalization}
Let $f:Y\to\PP^1$ satisfy the full CDP hypotheses
$b_1(Y)=b_2(Y)=0$ with smooth torus general fibre.
Let $mS$ be a full fibre with irreducible reduced support,
and let $Z\to T\to S$ be a resolution of its
normalization. If $Z$ is rational, then
\[
                              e(S)=2.
 \tag{C.MH.6}
\]
If $Z$ is ruled over a curve of genus $g\ge1$, then
\[
                              e(S)=2-2g.
 \tag{C.MH.7}
\]
The calculation retains all multiplicities in the full
fibre and places no normal crossings assumption on $S$.
\end{statement}
\begin{proof}
The integral class of D.1.9 restricts nontrivially to
$S$ and, by D.1.8, nontrivially to the nearby torus.
More precisely, restrict the section
$\gamma\in H^0(\PP^1,R^1f_*\mathbb Z)$ to a small
disc $\Delta$ about this fibre. The integral Leray edge
isomorphism
$H^1(f^{-1}\Delta,\mathbb Z)
 \simeq H^0(\Delta,R^1f_*\mathbb Z)$ supplies the
actual integral class on the tube: the two adjacent
base groups $H^1(\Delta,\mathbb Z)$ and
$H^2(\Delta,\mathbb Z)$ vanish.
Choose a prime $\ell$ for which its restriction to that
torus is nonzero modulo $\ell$. On a sufficiently small
fibre tube $U$, its reduction defines a connected cyclic
\'{e}tale cover with group $\mathbb Z/\ell$. Let $V$ be
the rational augmentation local system, that is, the
kernel of the trace from the regular representation
to the constant representation. Its rank is
$d=\ell-1$, it is self-dual, and $H^0(S,V)=0$.
Over $\mathbb C$ it is the sum of the nontrivial
characters of this cyclic group. Since $\ell$ is prime,
each of these characters restricts nontrivially to
the nearby torus $F$. The group-cochain complex for
$\mathbb Z^4$ with such a character is acyclic: one of
the four commuting operators $\rho(e_i)-1$ is
invertible, and its inverse gives a contraction of the
Koszul complex. Thus
\[
 H^j(F,V)=0\quad\hbox{for every }j,
 \qquad H^j(U\setminus S,V)=0\quad\hbox{for every }j.
 \tag{C.MH.8}
\]
The second assertion is the Leray sequence of the
punctured-disc torus bundle. This argument uses the
smooth fibres of the original map, so it also applies
when the full central fibre is $mS$.

Choose an embedded resolution of $(U,S)$. Its strict
transform is a blowup of $Z$, hence projective. Every
other exceptional surface is obtained by blowups from
a projective bundle over a compact smooth curve, or
from $\PP^2$. It is therefore projective. All centres
are compact and lie in the central fibre. C.2.MH.1
applies and gives
\[
                         H^3(S,V)\text{ pure of weight }3.
 \tag{C.MH.9}
\]

Let $A\subset T$ and $\Gamma\subset S$ be the reduced
conductor inverse image and image. The finite
normalization square is an actual topological pushout.
Its cohomological descent sequence, with the pulled
back local system, gives an exact sequence of mixed
Hodge structures
\[
\begin{gathered}
0\longrightarrow C_V\longrightarrow H^3(S,V)
  \longrightarrow H^3(T,V)\longrightarrow0,\\
C_V=\operatorname{coker}\bigl[H^2(T,V)\oplus H^2(\Gamma,V)
                         \longrightarrow H^2(A,V)\bigr].
\end{gathered}
 \tag{C.MH.10}
\]
This also covers the normal case, with the conductor
terms zero. The left term is pure of weight two:
$H^2(A,V)$ is the top cohomology of a compact curve
with a finite local system, and hence pure of weight
two. Strictness and (C.MH.9) make that term zero.
Consequently
\[
                           H^3(S,V)\simeq H^3(T,V).
 \tag{C.MH.11}
\]

For the resolution $\pi:Z\to T$, the restriction
$H^2(Z,V)\to H^2(E,V)$ to its exceptional curve is
surjective. Indeed $V$ is constant on each exceptional
fibre, since it is pulled back from $T$; the Gysin
classes of its irreducible exceptional curves restrict
by their negative-definite intersection matrix. That
matrix is invertible over $\mathbb Q$, separately
for each exceptional fibre and each coefficient
vector. The resolution descent sequence therefore
gives
\[
 H^3(T,V)\simeq H^3(Z,V),\qquad
 H^4(S,V)\simeq H^4(T,V)\simeq H^4(Z,V).
 \tag{C.MH.12}
\]
The last two identifications also follow directly
because all the omitted loci have real dimension two.

There is a second, purely topological consequence of
(C.MH.8). The support sequence and oriented local
Poincar\'e duality give
\[
 H^k(S,V)\simeq H_S^k(U,V)
             \simeq H_{6-k}(S,V).
 \tag{C.MH.13}
\]
Again these are the natural direct-limit maps on
sufficiently small tubes. Since $S$ has real
dimension four, (C.MH.13) gives $H^1(S,V)=0$ and,
using the self-duality of $V$,
$h^2(S,V)=h^4(S,V)$. Therefore (C.MH.11)--(C.MH.13)
give the exact numerical identity
\[
 d\,e(S)=\chi(S,V)=2h^4(Z,V)-h^3(Z,V).
 \tag{C.MH.14}
\]

If $Z$ is rational, the pullback of $V$ is constant,
$h^4(Z,V)=d$, and $h^3(Z,V)=0$; this proves (C.MH.6).
If $Z$ is ruled over $B$, blowups leave its fundamental
group and its third cohomology unchanged, and the
pulled-back local system comes from $B$. The
projective-bundle formula gives
$H^3(Z,V)=H^1(B,V)(-1)$. If the cyclic character is
trivial on $B$, its dimensions are
$h^4(Z,V)=d$ and $h^3(Z,V)=2gd$. Otherwise every
nontrivial character of the prime-order group remains
nontrivial on $B$, so $h^4(Z,V)=0$ and
$h^3(Z,V)=(2g-2)d$. Both substitutions in (C.MH.14)
give (C.MH.7).
\end{proof}


% C.2.MH.3 from part_C2_hodge.tex

